\documentclass[../main.tex]{subfiles}

\begin{document}

\section{Property of logarithms}

\paragraph{1} Reduce the following 

\section{Solving equations}

In this section are some solutions of different types of equations.

\begin{prob}{Example of exponential equation}{se-expfunc}

We are going to solve the following equation
\begin{gather*}
    8^{x-2} &= 3^{x+3}.
\end{gather*}

Since is an exponential function we apply a logarithmic function in order to apply logarithmic properties to manipulate the expression.


\begin{empheq}[box=\shadowbox]{equation*}
    \mathrm{Answer:}~(e^{2},1)
\end{empheq}
\end{prob}



\section{Modelling with exponential and logarithmic functions}

\subsection{Critical points}
To find the critical point we need to remember that for exponential functions we equate the argument to $0$ and then for $x$.
On the toher hand, for logarithmic functions we need to equate the argument of the function to \num{1} and solve for $x$.
Finally, we substitue the value of $x$ into the function and find the value of $f(x)$.

\begin{prob}{Example of exponential functions}{cp-expfunc}
Lets find the critical point of 
\begin{gather*}
    f(x) = \exp\qty[e\ln(x)-2e]
\end{gather*}

Since the argument is $e\ln(x)-2e$, lets equate to $0$ and solve for $x$.
\begin{align*}
    e\ln(x)-2e &= 0\\
    e\ln(x)-2e+2e &= 0+2e\\
    e\ln(x) &= 2e \\
    \frac{1}{e}e\ln(x) &= \frac{1}{e}2e \\
    \ln(x) &= 2 \\
    e^{\ln(x)} &= e^{2} \\
    x &=e^{2}
\end{align*}

Now that we know the value of $x$, we can substitute this value into the original equation to find the value of $f(e^2)$.
\begin{align*}
    f(x) &= \exp\qty[e\ln(x)-2e] \\
    f(e^{2}) &= \exp\qty[e\ln(e^{2})-2e] \\
        &=\exp\qty[2e\ln(e)-2e] \\
        &=\exp\qty[2e-2e] \\
        &=\exp\qty[0] = e^{0} \\
        &= 1.
\end{align*}
Hence the critical point is at $(e^{2},1)$.


\begin{empheq}[box=\shadowbox]{equation*}
    \mathrm{Critical~point~at:}~(e^{2},1)
\end{empheq}
\end{prob}


\begin{prob}{Example of logarithmic functions functions}{cp-logfunc}
Lets find the critical point of 
\begin{gather*}
        g(x) = \ln\qty[\exp[x^2]-8].
\end{gather*}

Since the argument is $\exp[x^2]-8=e^{x^2}$, lets equate to $1$ and solve for $x$.
\begin{align*}
    e^{x^2}-8 &= 1 \\
    e^{x^2}-8+8 &= 1+8 \\
    e^{x^2} &= 9 \\
    \ln\qty[e^{x^2}] &= \ln\qty[9] \\
    x^2\ln\qty[e] &= \ln\qty[9] \\
    x^2 &= \ln\qty[9] \\
    \qty(x^2)^{1/2} &= \qty(\ln\qty[9])^{1/2} \\
    x^{\frac{2}{2}} &= \qty(\ln\qty[9])^{1/2} \\
    x &= \sqrt{\ln\qty[9]}.
\end{align*}

Now that we know the value of $x$, we can substitute this value into the original equation to find the value of $f(\sqrt{\ln\qty[9]})$.
\begin{align*}
    g(x) &= \ln\qty[\exp[x^2]-8] \\
    f(\sqrt{\ln\qty[9]}) &= \ln\qty[\exp[\qty(\sqrt{\ln\qty[9]})^2]-8] \\
                         &= \ln\qty[\exp[\qty(\left(\ln\qty[9]\right)^{1/2})^2]-8] \\
                        &= \ln\qty[\exp[\qty(\left(\ln\qty[9]\right)^{2/2})]-8] \\
                        &= \ln\qty[\exp[\ln\qty[9]]-8] \\
                        &= \ln\qty[9-8] \\
                        &= \ln\qty[1] \\
                        &= 0
\end{align*}
Hence the critical point is at $(\sqrt{\ln\qty[9]},0)$.


\begin{empheq}[box=\shadowbox]{equation*}
    \mathrm{Critical~point~at:}~(\sqrt{\ln\qty[9]},0)
\end{empheq}
\end{prob}





\end{document}
